% ED26-BM-RAT2021-6-2: Over--Cruz's exact 2021 chapter, pp. 395--404.
\chapter{The Suppositional Theory of Conditionals and Rationality}
\label{sec:p19-rationality-6-2}

Over and Cruz's account concerns principally singular, natural-language
indicative conditionals: what happens on the next toss, or on this afternoon's
journey, rather than a universal rule about every member of a domain.
Its central proposal is that assessing ``if \(A\), then \(B\)'' amounts to
assessing \(B\) while supposing \(A\). The associated degree of belief is
the subjective conditional probability \(P(B\mid A)\). The authors propose
both a normative standard and a descriptive hypothesis. Probability-theoretic
coherence supplies the former; correspondence between conditional judgments
and conditional probabilities supplies the latter. A descriptive fit does not
establish the normative premise, and endorsing that premise does not entail
that people always comply with it \cite[pp.~395--401]{OverCruz2021Suppositional}.

\section{The account and the prediction to be tested}
\label{sec:p19-r62-account}

\begin{hbdefinition}{def:p19-r62-equation}{The suppositional account and its two claims}
Write \(P(\text{if }A\text{ then }B)\) for the subject's assessment of
the ordinary-language conditional, without presupposing that it is the
probability of a fixed Boolean event. The \emph{Equation} is
\[
 P(\text{if }A\text{ then }B)=P(B\mid A).
\]
The chapter's normative proposal assesses reasoning with that interpretation
by coherent subjective probability. Its \emph{conditional probability
hypothesis} predicts that people's judgments of the conditional will
correspond to their judgments of \(P(B\mid A)\). This is a prediction
about a specified elicitation task, not a definition of the observed response.
For the finite calculations below, ratio conditioning requires \(P(A)>0\).
The psychological operation of supposing an antecedent is not defined by
division by zero when that condition fails.
\end{hbdefinition}

The historical foundations invoked are Ramsey's 1926 and 1929 writings,
de Finetti's 1936 and 1937 work, and Edgington's development of the Equation.
The chapter also calls its conditional a probability conditional or
conditional event. Those names identify an interpretation here, not an
automatic adoption of every axiom for iterated conditional random quantities
in \cref{sec:p15-rationality-4-4}.

Reasoning from \(A\), together with background beliefs, to \(B\) can
support a conditional assertion; the assertion can in turn summarize that
reasoning. The authors' unification proposal therefore places conditional
probability at the interface between reasoning, judgment and decision making.
It is a research program, not a completed theorem equating every cognitive
task or every aspect of rationality. Cognitive limitations and fallacies are
compatible with the descriptive hypothesis. The authors predict that
probability theory will describe conditional reasoning more accurately
than binary extensional logic. That comparative proposal needs matched
tasks and model versions; it is not established for every task by the
Equation's definition.

The source reports corroboration across Evans--Handley--Over (2003),
Evans--Handley--Neilens--Over (2007), Oberauer--Wilhelm (2003),
Over and colleagues (2007), Pfeifer--Kleiter (2005), Fugard and colleagues
(2011), and Singmann--Klauer--Over (2014). Its review does not provide those
studies' complete recruitment, exclusions, response distributions or
uncertainty estimates. The original-study records in
\cref{sec:p19-rat46-interpretations,sec:p19-r45-coherence-studies} retain
their own tasks and denominators; they are not pooled into a population
frequency by this list of corroborating citations.

\begin{hbexample}{def:p19-r62-singular-general}{A bird rule is not a next-animal prediction}
The assertion that birds fly, read as a general conditional covering its
whole domain, has counterexamples. A prediction that the next animal seen
will fly if it is a bird can nevertheless be highly credible in a suitable
setting. Its credibility depends on the local distribution and background
knowledge. The chapter confines its main treatment to singular conditionals,
referring the comparison of ``if'' with ``all'' to Cruz--Oberauer (2014).
Success on one task does not establish a semantics for the other.
\end{hbexample}

\section{Material implication, the coins, and the scope of the objection}
\label{sec:p19-r62-material}

\begin{hbconvention}{def:p19-r62-values}{Four cases and four assessments}
In the following array, \(1,0\) mean factual truth and falsity in the
first four columns; \(V\) means void in the de Finetti column. In the
Jeffrey column, the entries on false-antecedent cases are subjective values,
not a third factual truth value or the number of observed responses.
Put \(r=P(B\mid A)\), when this ratio is defined.
\[
\begin{array}{cc|c|c|c|c}
 A&B&A\lor B&\neg A\lor B&\text{de Finetti if}&\text{Jeffrey if}\\\hline
 1&1&1&1&1&1\\
 1&0&1&0&0&0\\
 0&1&1&1&V&r\\
 0&0&0&1&V&r
\end{array}
\]
The disjunction and material columns are ordinary bivalent truth tables.
The last two reproduce the distinct value assignments of the source's
Tables~6.2.3--6.2.4. Their interpretation must not be collapsed into
material implication merely because they agree on antecedent-true cases.
\end{hbconvention}

The next toss of a fair coin is heads with conditional probability \(.5\).
The next toss of a double-headed coin is tails with conditional probability
zero. These assessments do not become high merely because the two ancient
coins are locked in a museum and unlikely to be tossed. Material implication
does become highly probable in that circumstance, because every
false-antecedent case satisfies \(\neg A\lor B\).

\begin{hbtheorem}{thm:p19-r62-material}{proved=annex:p19.r62.material}{The material prediction and its exact coincidence cases}
On a finite probability space, let \(a=P(A)>0\) and
\(r=P(B\mid A)\). Then
\[
 P(\neg A\lor B)=1-a+ar,\qquad
 P(\neg A\lor B)-r=(1-a)(1-r).
\]
The two assessments agree exactly when \(a=1\) or \(r=1\).
If all four \(A,B\) cases have mass \(1/4\), the conditional-probability
assessment is \(1/2\), the material assessment \(3/4\).
If the fair-coin toss has probability \(a>0\), the material assessment
is \(1-a/2\); for a double-headed-coin toss of probability \(a>0\),
the tails material assessment is \(1-a\). Both approach one as tossing
becomes unlikely, although their corresponding conditional probabilities
remain \(1/2\) and zero.
\end{hbtheorem}

Or-introduction preserves probability: since \(A\) implies
\(A\lor B\), coherence requires \(P(A)\leq P(A\lor B)\).
Consequently identifying a natural-language conditional with its material
counterpart makes it follow from \(\neg A\) and forces its probability
to be at least \(P(\neg A)\). The museum example exposes the undesirable
prediction under the proposed suppositional interpretation. It is not a
countermodel to the classical theorem about material implication.
Over and Cruz report rejection of this natural-language paradox in
Cruz--Over--Oaksford (2017), Pfeifer--Kleiter (2011), and
Politzer--Baratgin (2016); rejection of a conversational inference and
rejection of its probability inequality remain distinct responses.

The chapter's comparator is the 1991 mental-model account in which the
paradoxes were valid. It explicitly acknowledges the 2015 revision by
Johnson-Laird, Khemlani and Goodwin, in which those paradoxes are no longer
valid, and declines to assess that developing revision here. The selected
2021 mental-model treatment at \cref{sec:p19-rationality-2-3} therefore
cannot be assigned every prediction of the older comparator by its name.

The chapter further contrasts Lewis-style closest-world conditionals with
suppositional ones. A closest-world conditional can be evaluated at a world
where its antecedent is false; the distribution over those worlds can affect
its probability, whereas ratio \(P(B\mid A)\) depends on the
antecedent-true portion. It invokes Lewis's 1976 general obstruction to
identifying those quantities. That imported result is not proved by the
museum example. The finite fixed-Boolean-event obstruction at
\cref{thm:p19-rat46-fixed-event} is a separately proved restricted bridge,
not the full Lewis theorem or a prohibition of every conditional semantics.

\section{Supposition, factual evaluation, and bets}
\label{sec:p19-r62-supposition}

\begin{hbmethod}{def:p19-r62-ramsey}{The Ramsey-test proposal}
For a conditional assessment, hypothetically suppose the antecedent, make
the changes needed to preserve or restore consistency with background
beliefs, and assess confidence in the consequent within that supposition.
The source describes minimal changes. This specifies a psychological
operation to investigate; it does not select a unique revision operator,
distance on belief states, or numerical policy for every impossible or
probability-zero supposition. The closest-world semantics of
\cref{def:p08-stalnaker-conditional} and this judgment procedure must be
identified separately before transferring a theorem between them.
\end{hbmethod}

For the source's afternoon journey, the short city route \(S\) may be
quick \(Q\), and the longer countryside route \(L\) enjoyable \(E\).
Supposing each route calls on beliefs about traffic, weather and accidents.
If both conditional prospects are good, choosing between them still depends
on the relative value of speed and enjoyment. Conditional credibility alone
does not determine utility or a uniquely rational route.

After taking the countryside route, the city-route indicative is no longer
operative for the actual journey. The chapter follows Ramsey and de Finetti
in calling it void in that respect. One can instead consider the
counterfactual that the city journey would have been quick. It connects
that counterfactual assessment with the regret illustration attributed to
Kahneman--Miller (1986): an accident on an unusual route can elicit more
regret than an accident on the usual route. That is a source-reported
counterfactual/regret result, not a new measurement of drivers or proof that
regret improves choices. The chapter notes both rational and irrational
effects on future decisions, then returns to indicative conditionals.

The de Finetti column distinguishes confirmed, refuted and void cases.
Calling a traditional truth-table response ``defective'' names the departure
from a bivalent material table; Over and Cruz argue that it instead fits
this conditional assessment, referring to Baratgin--Over--Politzer (2013),
Politzer--Over--Baratgin (2010), and Over--Baratgin (2017). The third
response is not evidence of defective reasoning by definition.
Ramsey's exception for consequences of laws or hypotheses also matters:
the account treats a conditional from \(A\land B\) to \(B\) as
certain, rather than simply discarding it because its antecedent is false.
This is a commitment about logical supposition, not the arithmetic claim
that \(0/0=1\).

The Jeffrey column retains a subjective conditional assessment on a
false-antecedent case while withholding a factual verdict. The following
calculation explains its probability interpretation without turning that
value into bivalent truth.

\begin{hbtheorem}{thm:p19-r62-jeffrey-value}{proved=annex:p19.r62.jeffrey-value}{The Jeffrey-table random quantity and the null-antecedent boundary}
For \(r\in[0,1]\), define the finite random quantity \(X_r\) to be
one on \(A\land B\), zero on \(A\land\neg B\), and \(r\)
on \(\neg A\). If \(P(A)>0\),
\[
 E_P(X_r)=r\quad\text{if and only if}\quad r=P(B\mid A).
\]
If \(P(A)=0\), the same expectation equation holds for every
\(r\in[0,1]\); it determines no unique conditional assessment.
This is an explicit finite representation of the chapter's table,
not a claim that participants compute expectations or a complete
coherence theorem for arbitrary families of conditional bets.
\end{hbtheorem}

\section{Uncertain premises and probabilistic validity}
\label{sec:p19-r62-coherence}

\begin{hbdefinition}{def:p19-r62-p-valid}{The source's validity criterion}
For an assessed statement with probability \(p\), its uncertainty is
\(1-p\). A one-premise inference is \emph{p-valid} when every coherent
assignment makes the conclusion at least as probable as its premise.
For finitely many premises, p-validity requires, under every coherent
assignment allowed by the stated interpretation, that conclusion
uncertainty be no greater than the sum of premise uncertainties.
This definition, attributed to Adams (1998), is not a significance
test and is not the rate at which participants endorse the inference.
A strong inductive argument can increase uncertainty and is not thereby
a p-valid deduction.
\end{hbdefinition}

Conjunction elimination illustrates the criterion without an uncertain
conditional premise: \(A\land B\) entails \(B\), and coherent
confidence in \(B\) cannot be below confidence in \(A\land B\).
The conjunction fallacy is the reverse numerical ordering under this
event interpretation. For a conditional premise, the relevant model is
the suppositional one, not the material column in
\cref{def:p19-r62-values}.

\begin{hbtheorem}{thm:p19-r62-intervals}{proved=annex:p19.r62.intervals}{Elimination, centering, and modus ponens under declared assessments}
On the full four-atom algebra for \(A,B\), with freely varying
nonnegative masses, the following intervals are sharp.
Given \(P(A\land B)=x\), the allowed values of \(P(B)\) form
\([x,1]\). If additionally \(x>0\), the allowed values of
\(P(B\mid A)\) also form \([x,1]\).
Given \(P(A)=a>0\) and \(P(B)=b\), the allowed values of
\(P(B\mid A)\) form
\[
 \left[\max\left(0,\frac{a+b-1}{a}\right),
       \min\left(1,\frac ba\right)\right].
\]
Both one-premise and two-premise centering are p-valid on these
positive-antecedent assessments. Given \(P(A)=a>0\) and
\(P(B\mid A)=r\), the allowed \(P(B)\) values form
\([ar,ar+1-a]\), and modus ponens is p-valid.
The last interval gives \([.4,.9]\) for \(a=.5,r=.8\), and
the singleton \(\{.8\}\) for \(a=1,r=.8\).
These are finite coherence statements, not fitted cognitive algorithms.
\end{hbtheorem}

The countryside example has a high but not certain conditional assessment,
because weather and other disabling conditions may make the journey
unpleasant. Knowing that the chosen route is \(L\) supports confidence
\(.8\) in enjoyment when \(P(E\mid L)=.8\); uncertainty about
whether \(L\) is chosen broadens the interval to \([.4,.9]\) at
\(P(L)=.5\). An unspecified \(P(E\mid\neg L)\) does not justify
selecting its midpoint or treating the conclusion as certain.

The normative interpretation is conditional on probability coherence and
the declared conditional interpretation. Over and Cruz connect incoherence
with exposure to a Dutch book and cite Howson--Urbach (2005) for the Bayesian
justification. The finite stake-based coherence construction and its
null-stratum qualifications are developed at \cref{sec:p15-rationality-4-4};
this chapter's interval calculation does not establish an unconditional
psychological obligation to make every available bet. Nor does coherence
alone ensure accurate priors, good utilities, or an optimal use of effort.

\subsection{Does making an inference improve coherence?}
\label{sec:p19-r62-explicit-inference}

The chapter gives a reason to test explicit inference: if it costs time and
effort, it might improve the coherence of the resulting beliefs. It reports
improvements in some studies but not a general benefit, citing Cruz and
colleagues (2015, 2016, 2017), Evans--Thompson--Over (2015), and
Singmann--Klauer--Over (2014). In particular, it emphasizes the persistence
of the conjunction fallacy even when the constituent proposition is
explicitly inferred from the conjunction. Prototypicality and content
matter; an inference form alone does not determine response quality.
The chapter points to Costello--Watts (2014) and Tentori--Crupi--Russo
(2013) as contrasting accounts relevant to explaining the conjunction
fallacy. The former's probability-plus-noise account and the latter's
probability/inductive-confirmation distinction are not empirical
consequences of p-validity. The persistence of the fallacy constrains
an explanation; the chapter does not adjudicate those explanations
with a new experiment.

An inspected original anchor is Cruz and colleagues' two-experiment
2015 article. Experiment~1 retained 871 of 1,140 completed online records
from English-speaking countries after duplicate-IP, invariant-response and
age exclusions. Respondents used percentage confidence for separate
statements or explicit premise--conclusion pairs, with Linda and contrasting
stereotype contents. Weighted random assignment to differently numerous
booklets yielded 305 statement and 566 inference participants. The retained
sample's mean age was 35, range 12--78; twenty were under eighteen, and
the authors report the same significance pattern after excluding them.
It was a CrowdFlower sample, not probability sampling of all reasoners.

The measure subtracted the coherence interval's width from a binary
compliance score. Its chance baseline assumes a uniform random judgment
over the probability scale. The original article calls this an approximation
and notes possible boundary biases. Overall above-chance coherence was
reported as \(F(1,869)=885.29,p<.001\), but no significant task benefit
was found. That result is not an equivalence test and does not show that
explicit inference never helps. Different interval widths also make
above-chance scores across forms different from raw compliance rates
\cite[pp.~3--5,8]{CruzEtAl2015Coherence}.

Experiment~2 used 48 volunteering Orsay students, native French speakers,
mean age 20, range 18--24, mainly biology or medicine students, in two
equal statement/inference groups. They assessed conjunction--conditional
and conjunction-elimination forms with prototypicality varied, at their
own pace in the university library. Centering judgments were above the
authors' chance baseline, \(F(1,46)=36.83,p<.001\), without a
significant task effect. The conjunction-elimination content contrast was
\(F(1,46)=33.31,p<.001\); counter-prototypical judgments were below
that baseline, \(F(1,46)=18.37,p<.001\), even with explicit inference.
Those statistics are not participant proportions or population prevalence.
The corrigendum reverses the original Figure~2's third and fourth plotted
positions. Its 95\% interval bars do not print exact endpoint values;
no digitized endpoints or new effect intervals are asserted here
\cite[pp.~5--8]{CruzEtAl2015Coherence}
\cite{CruzEtAl2015CoherenceCorrigendum}.

\section{Centering, missing links, and a rival semantics}
\label{sec:p19-r62-centering}

One-premise centering moves from \(A\land B\) to a conditional from
\(A\) to \(B\); two-premise centering starts with the two conjuncts
separately. Their probabilistic bounds were established in
\cref{thm:p19-r62-intervals}. The source also accepts centering for
material and centered closest-world interpretations, but contrasts
inferentialist accounts developed by Douven (2015, 2016) and
Douven--Elqayam--Singmann--van Wijnbergen-Huitink (2018).
In the semantic inferentialist version described, the conditional requires
a deductive, inductive or abductive connection from antecedent to consequent.
The truth of both components need not establish that connection.

\begin{hbexample}{def:p19-r62-missing-link}{Sun in Canberra and rain in London}
Let \(C\) denote sunshine in Canberra today and \(N\) rain in London.
The chapter asks whether knowing \(C\) and \(N\) warrants the conditional
from \(C\) to \(N\). The stipulated suppositional standard supports
centering; semantic inferentialism can reject it for lack of a link.
But opening a conversation with the two facts and then drawing that
conditional can be pragmatically peculiar. Conversely, presenting the
conditional first and then confirming both facts asks for a verdict on
an already-presented assertion, not its usefulness as a newly inferred
conclusion. A conditional bet confirmed on an antecedent-true,
consequent-true outcome raises a further operational question.
These task descriptions are not interchangeable observations.
\end{hbexample}

Over and Cruz cite Skovgaard-Olsen--Kellen--Krahl--Klauer (2017,
Figure~1) for affirmative truth judgments in the conditional-first task.
They invoke earlier truth-table work for the same confirmed-case verdict,
and Cruz and colleagues (2015, 2016), Pfeifer--Tulkki (2017), and
Politzer--Baratgin (2016) for centering evidence. The 2021 chapter supplies
neither the full 2017 sampling/design record nor a common effect estimate
across these tasks. Its asserted task-order distinction should be tested
with task-specific original records, not treated as an experimentally
isolated universal order effect.

If \(A,B\) are probabilistically independent, \(P(B\mid A)=P(B)\)
for a positive antecedent: observing \(A\) provides no probabilistic
information about \(B\). The chapter takes this to explain why a
missing-link conditional need not be useful in a decision. Neither
absence of a conversational connection nor different geographical domains
alone establishes probabilistic independence; and an evidential
conditional is not automatically an intervention prediction about doing
\(A\). The claimed lack of decision usefulness therefore does not
follow from syntax alone.

For relevant conditionals, Skovgaard-Olsen--Singmann--Klauer (2016) is
reported as finding better Equation conformity when \(P(B\mid A)>
P(B\mid\neg A)\). The chapter also cites counterevidence from
Oberauer--Weidenfeld--Fischer (2007), Over and colleagues (2007), and
Singmann and colleagues (2014). These are different tasks and materials,
not replications with a single shared denominator. The separately qualified
2016 relevance record at \cref{sec:p19-rat46-relevance} retains its
accepted-manuscript version and exclusions.

The authors' central response is the report that Cruz and colleagues (2016)
found missing-link effects on conjunction and disjunction judgments too.
An effect shared with other connectives cannot, by itself, identify what
distinguishes the semantics of conditionals. They infer a pragmatic
explanation, emphasizing a shared topic of discourse rather than a uniquely
conditional epistemic connection. They also mention Sanborn--Chater (2016)
for the possible difficulty of combining probabilities across domains.
The evidence defeats an inference from a non-specific effect to a
conditional-specific meaning; it does not prove that semantics never
contributes, or that every missing-link rejection is pragmatic.
\Cref{sec:p19-rat46-relevance} preserves Oberauer and Pessach's less
settled assessment of the semantics/pragmatics dispute.

\section{Belief updating and dynamic modus ponens}
\label{sec:p19-r62-updating}

The chapter's scientific example starts with uncertain belief in a
hypothesis \(H\) but high confidence in \(H\) supposing an experimental
finding \(F\). Learning \(F\), and nothing else relevant, can then
justify changing belief in \(H\). It calls the resulting use of the
conditional \emph{dynamic modus ponens}. This is an updating claim over
time, not a simultaneous assignment of two conflicting probabilities to
the same event under one fixed law.

\begin{hbdefinition}{def:p19-r62-invariance}{The premise needed for the source's update}
Let \(P_1\) be a finite prior with \(0<P_1(F)<1\), and
\(t=P_2(F)\) the assessed new marginal. Jeffrey updating stipulates
preservation of the prior conditional laws within \(F\) and
\(\neg F\). In particular, for an event \(H\), put
\(r=P_1(H\mid F)\), \(q=P_1(H\mid\neg F)\), and use
\[
 P_2(H)=tr+(1-t)q.
\]
For a new law with both strata positive, preservation means
\(P_2(H\mid F)=r\) and \(P_2(H\mid\neg F)=q\).
An endpoint update has a null stratum and does not define its conditional
law by a ratio. Preservation is a substantive assumption about the
information received, not a consequence of merely learning a new marginal.
\end{hbdefinition}

\begin{hbtheorem}{thm:p19-r62-update}{proved=annex:p19.r62.update}{The source's updating calculation and what it does not determine}
Under \cref{def:p19-r62-invariance}, the mixture of the two prior
conditional laws is a probability law with marginal \(t\) for \(F\).
At \(t=1\), it gives \(P_2(H)=r\). If only \(r,t\) are fixed
and the prior false-finding conditional \(q\) is unconstrained beyond
\([0,1]\), it permits exactly \([tr,tr+1-t]\) for \(P_2(H)\).
Thus \(r=.8,t=.7\) gives \([.56,.86]\) under those limited inputs.
An additional fixed prior marginal for \(H\) or \(F\) can narrow
that interval; it is not automatically sharp relative to those further
constraints. Without preservation, a new marginal \(t\) alone permits
every new \(P_2(H)\in[0,1]\) on a full four-atom algebra.
\end{hbtheorem}

Over and Cruz's numerical example begins with \(P_1(H)=.5\) and
\(P_1(H\mid F)=.8\). Certain learning of \(F\) gives \(.8\)
under the stipulated update. Their less-certain \(P_2(F)=.7\)
requires the false-finding conditional as well. The prior marginal
\(.5\) does not specify it without more information about \(P_1(F)\).
For a fully specified \emph{illustrative bridge}, choosing
\(P_1(F)=.5\) makes \(q=.2\), with four prior masses
\((FH,F\neg H,\neg FH,\neg F\neg H)=(.4,.1,.1,.4)\).
The invariant update yields \(P_2(H)=.7(.8)+.3(.2)=.62\).
Those extra masses are stipulated here, not measured participants' beliefs
or additional values given by the 2021 chapter.

The source cites Hadjichristidis--Sloman--Over (2014) and Zhao--Osherson
(2010) for failures to conform perfectly to Jeffrey updating. It proposes,
as a possible explanation, concentrating on the more probable weighted
term and neglecting the other. This is a hypothesis about process, not
identified by a discrepancy from the formula alone. The inspected Zhao
study record at \cref{sec:p19-r45-updating-studies} distinguishes the
flashlight/color evidence task, additional certainty comparison, and
new narrated lottery/game sample. The legitimacy of invariance differs
with the information mechanism. Changing a conditional link can be an
appropriate revision, not necessarily a probability fallacy.

\section{Interfaces and limits of the proposed unification}
\label{sec:p19-r62-interfaces}

\Cref{sec:p15-rationality-4-4} develops the coherence calculus, conditional
events and possible-null strata; it does not infer a cognitive mechanism
from its propagation theorems. \Cref{sec:p19-rationality-4-5} gives
Chater and Oaksford's knowledge-rich Bayesian program, with explicit
observation and updating assumptions. \Cref{sec:p19-rationality-4-6}
distinguishes probability-of-conditional, truth, acceptance and inference
tasks, and assesses competing response models more cautiously. These
source-specific claims support comparison, not a single amalgamated account.

Suppression in \cref{sec:p19-conditional-reasoning-suppression} and
possibility-premise effects in
\cref{sec:p19-conditional-interpretation-experiments} concern endorsements
under particular instructions. A changed background distribution, a
changed interpretation, and an unchanged material inference with lower
endorsement are distinct explanatory possibilities. Neither the Equation
nor p-validity alone supplies the response process for those experiments.
Likewise a dual-process description at \cref{sec:p19-rationality-2-5}
cannot be inferred from coherence by relabeling correct judgments as one
process and errors as the other.

The source's conclusion is that suppositional conditionals and coherence
offer a framework for investigating both reasoning and decision making,
including rational and irrational belief updating. It leaves the extent
of dynamic coherence and the invariance of subjective conditional
probabilities as future empirical questions. The worked success case is
an explicitly invariant update with a declared joint model; failure cases
include treating a rare antecedent as sufficient for high conditional
belief, deriving one updating value from an unspecified information
mechanism, or interpreting a numerical/task response as a normative
verdict without fixing its meanings. A broader account remains a proposal
to be tested, not a finding that people are uniformly Bayesian or
uniformly incoherent \cite[pp.~399--401]{OverCruz2021Suppositional}.
